\scanpage{296}
\mspage{- 2 -}

\uline{il existe une suite spectrale de} $\handwritten{\Lambda}$\uline{-modules aboutissant à} $\mathrm{Ext}_{\handwritten{\Lambda}}^{\handwritten{*}}(\underset{\handwritten{\longrightarrow}}{\lim}\handwritten{A}_{\handwritten{\alpha}},B)$, \uline{et dont le premier terme est donné par}

\[ E_2^{pq}=\underset{\handwritten{\longleftarrow}}{\lim}_{\handwritten{\alpha}}^{(p)}\mathrm{Ext}_{\handwritten{\Lambda}}^{q}(A_{\handwritten{\alpha}},B). \]

Utilisons maintenant cette proposition avec $B=\handwritten{\Lambda}$ (régulier) et $A_{\handwritten{\alpha}}=\mathrm{Ext}_{\handwritten{\Lambda}}^{k}(L_k(M_{\handwritten{\alpha}}),\handwritten{\Lambda})$. Nous obtenons $E_2^{pq}=0$, $q\handwritten{<}k$; $E_2^{pk}=\underset{\handwritten{\longleftarrow}}{\lim}^{(p)}L_k(L_k(M_{\handwritten{\alpha}}))=\underset{\handwritten{\longleftarrow}}{\lim}^{(p)}L_k(M_{\handwritten{\alpha}})$ et $E_2^{pq}=0$, $p\handwritten{>}\dim \handwritten{\Lambda}-q-2$, $q\handwritten{>}k$, en vertu du théorème 1, resp. $\handwritten{\alpha})$, $\handwritten{\eta})$, $\handwritten{\vartheta})$ de ($^{1}$) et à cause de l'hypothèse de récurrence. Donc la suite spectrale dégénère en l'isomorphisme

\[ \underset{\handwritten{\longleftarrow}}{\lim}^{(p)}L_k(M_{\handwritten{\alpha}})=\mathrm{Ext}_{\handwritten{\Lambda}}^{p+k}(\underset{\handwritten{\longrightarrow}}{\lim}\mathrm{Ext}_{\handwritten{\Lambda}}^{k}(L_k(M_{\handwritten{\alpha}}),\handwritten{\Lambda}),\handwritten{\Lambda})\quad\text{pour }p\handwritten{\geq}\dim \handwritten{\Lambda}-k, \]

et le \deleted{\unsure{de|le}} deuxième membre est bien zéro pour $p\handwritten{>}\dim \handwritten{\Lambda}-k$, ce qui achève l'induction et la démonstration du théorème 2.

REMARQUE 1. - Le théorème 2, restreint au cas $\handwritten{\Lambda}=\handwritten{\mathbb{Z}}$, généralise un résultat de Milnor (($^{3}$), Appendix).

REMARQUE 2. - Le théorème 2 n'est pas valable en général si $I$ n'est pas filtrant, ou si les $M_{\handwritten{\alpha}}$ ne sont pas de type fini.

\margin{\handwritten{$\checkmark$}}

En ce qui concerne les $\underset{\handwritten{\longleftarrow}}{\lim}^{(i)}M_{\handwritten{\alpha}}$ avec $i\handwritten{\leq}\dim \handwritten{\Lambda}$, nous n'avons que des résultats fragmentaires. A l'aide des méthodes de ($^{2}$), on démontre facilement

% slanted
\margin{\handwritten{A \uline{\unsure{signaler|vérifier}} ?}}

PROPOSITION 2. - \uline{Soit} $A$ \uline{un} $\handwritten{\Lambda}$-\uline{module} ($\handwritten{\Lambda}$ commutatif) \uline{et} $(M_{\handwritten{\alpha}})$ \uline{un système projectif de} $\handwritten{\Lambda}$-\uline{modules. Alors il existe deux suites spectrales} de $\handwritten{\Lambda}$-modules \uline{aboutissant à la m\handwritten{\^}{e}me limite} $H^{\handwritten{*}}$, \uline{avec comme premiers termes}

\[ E_2^{pq}=\underset{\handwritten{\longleftarrow}}{\lim}^{(p)}\mathrm{Ext}_{\handwritten{\Lambda}}^{q}(A,M_{\handwritten{\alpha}})\quad\text{\uline{et}}\quad {}'E_2^{pq}=\mathrm{Ext}_{\handwritten{\Lambda}}^{p}(A,\underset{\handwritten{\longleftarrow}}{\lim}^{(q)}M_{\handwritten{\alpha}}). \]

Soit maintenant $\handwritten{\Lambda}$ régulier, $I$ filtrant et les $M_{\handwritten{\alpha}}$ et $A$ de type fini. Alors $E_2^{pq}=0$ pour $p\handwritten{>}\dim \handwritten{\Lambda}-q$ $\handwritten{[}$théorème 1, $\handwritten{\delta})$ de ($^{1}$) et théorème 2$\handwritten{]}$, d'où $H^s=0$ pour $s\handwritten{>}\dim \handwritten{\Lambda}$. La deuxième suite spectrale donne $\mathrm{Ext}_{\handwritten{\Lambda}}^{1}(A,\underset{\handwritten{\longleftarrow}}{\lim}^{(1)}M_{\handwritten{\alpha}})=0$, si $\dim \handwritten{\Lambda}=1$, et $\mathrm{Ext}_{\handwritten{\Lambda}}^{\handwritten{\nu}-1}(A,\underset{\handwritten{\longleftarrow}}{\lim}^{\handwritten{(\nu)}}M_{\handwritten{\alpha}})\handwritten{=}=0$, si $\handwritten{\nu}=\dim \handwritten{\Lambda}\handwritten{>}1$, pour tout $\handwritten{\Lambda}$-module $A$ de \uline{type fini}. Mais ceci

% archive sheet number: 295

